\section{Analytic angular germs and nonlinear birth of exact labels}
\label{sec:v68-angular-germs}
A zero tangent fraction does not imply a zero physical source. We
resolve this distinction for the original circular table. The exact
angular sets have an analytic one-sided germ, even when several active
boundaries have the same tangent. A nonzero germ consequently has a
finite order of birth. That order is not assumed bounded as the word
length increases.

Throughout Sections~\ref{sec:v68-angular-germs}--\ref{sec:v68-annular-height},
fix $0<\chi\le1/10$ and $0<\varepsilon\le\chi/4$. Use the
selected-contact charts of Theorem~\ref{thm:v66-uniform-morse}, their
original physical restrictions, and the exact label $\lambda=(n,k)$.
The radius $r_\chi$, roof width $h_\chi=r_\chi^2/2$, density
constant $K_\chi$, and coefficient $J_z$ retain their previous meanings.
The selected orbit is continued only to define coordinates and masks;
no continued nonphysical point is assigned physical measure.

\subsection{Analyticity of the primitive tests}
\begin{lemma}[Analytic physical decision germs]
\label{lem:v68-analytic-margins}
At a fixed normal critical word with positive selected incidences, the
complete finite scalar-margin description of
Lemma~\ref{lem:v65-margin-description} can be chosen real analytic in
a neighborhood of the origin in the Morse chart. Every active margin
has nonzero differential. The inactive signs remain fixed on a smaller
positive radius, which is not asserted uniform in the word.
\end{lemma}
\begin{proof}
For the circular table, the selected contact equations are analytic
functions of contact arclengths and the two endpoint variables. The
flight lengths are positive. Their internal derivative is invertible
by Lemma~\ref{lem:v66-weighted-inverse}, so the analytic implicit
function theorem gives the selected contact graph at this fixed word.
The coordinate construction in the proof of
Theorem~\ref{thm:v66-uniform-morse} is analytic as well: its matrix
$Q_z(x)$ is analytic and positive definite, matrix square root is analytic
on the positive definite cone, and the map $x\mapsto Q_z(x)^{1/2}x$
has an analytic local inverse. We keep this specific Morse map.

The original section is a finite union of rectangles in the analytic
$(\alpha,p)$ charts; see \eqref{eq:enlarged-section}. Its edge tests,
composed with each regular selected branch, are analytic. On a
near-clearance envelope the perpendicular foot lies strictly inside
the flight and the signed perpendicular distance is nonzero at the
center. Its sign can therefore be fixed before taking its absolute
value. The resulting clearance margin is analytic. The completeness
and the nonzero differentials are those of
Lemma~\ref{lem:v65-margin-description}; all other first-hit tests are
retained as strict inactive tests on the chosen disk. Taking a minimum
is needed for the smooth guard, but not for the scalar sign list: its
individual candidates are kept separately. The explicit inactive-sign
argument in \eqref{eq:v67-inactive-radius} gives the last assertion.
No claim of a count-uniform analytic radius or derivative bound is made.
\end{proof}

Let $B_{z,\lambda}(y)$ be the complete physical and exact-label
indicator in \eqref{eq:v66-exact-label-partition}. For $r>0$ put
\begin{equation}\label{eq:v68-unweighted-angular-germ}
 \ell_{z,\lambda}(r)=\frac1{2\pi}\int_0^{2\pi}
                       B_{z,\lambda}(r\omega_\phi)\,d\phi.
\end{equation}
This is the unweighted angular fraction. It is distinct from the
clearance-guard-weighted angular fraction in the next section and
from the roof density. The initial section test, occupation at
$0,\ldots,m-1$, terminal membership at $m$, and displacement are
all evaluated before a label is assigned.

\begin{lemma}[Ordering the analytic angular boundaries]
\label{lem:v68-root-order}
For each fixed word there is $r_z>0$ on which every active scalar
margin has exactly two angular zero germs. They are real analytic
functions of $r$ at zero. After coincident germs have been identified,
their cyclic order is fixed for $0<r<r_z$. The label on each open
arc between consecutive germs is independent of $r$.
\end{lemma}
\begin{proof}
For an active analytic margin $f$, the function
\[
 F_f(r,\phi)=\begin{cases}f(r\omega_\phi)/r,&r\ne0,\\
                           Df(0)\cdot\omega_\phi,&r=0
              \end{cases}
\]
is analytic. At either zero $\phi_0$ of its value at $r=0$,
\[
 \partial_\phi F_f(0,\phi_0)=Df(0)\cdot\omega'_{\phi_0}\ne0.
\]
The analytic implicit function theorem supplies one root near each
of these two angles. On the compact complement of their neighborhoods,
$F_f(0,\cdot)$ stays away from zero, so there are no further roots
for small $r$. Choose an angular cut avoiding every zero at $r=0$.
All roots then have analytic lifts to one interval of length $2\pi$.

The difference of two such lifts is either the zero analytic germ or
has a first nonzero Taylor coefficient. In the latter case it has a
fixed sign for all sufficiently small positive $r$. There are finitely
many pairs. Shrinking the radius once fixes their order simultaneously.
Identifying a coincident root does not remove any primitive test: all
of the sign changes attached to that root are retained. On each open
moving arc no active test vanishes, and all inactive signs are fixed.
The full input sign vector is therefore constant there as $r$ varies.
The original one-hot visit rule assigns the same label, or no label,
to that arc. Boundary conventions concern only finitely many angles
at each positive radius.
\end{proof}

\begin{theorem}[Complete finite-order classification of angular germs]
\label{thm:v68-angular-germ-classification}
For every fixed selected word and exact label, the function
$\ell_{z,\lambda}$ has a real analytic extension from a sufficiently
small right neighborhood of zero. Its value at zero is the tangent
fraction $\theta_{z,\lambda}$. Precisely one of the following holds:
\begin{equation}\label{eq:v68-germ-alternatives}
 \ell_{z,\lambda}\equiv0\quad\hbox{near zero},\qquad\hbox{or}\qquad
 \ell_{z,\lambda}(r)=\eta_{z,\lambda}r^{d_{z,\lambda}}
                              +O_z(r^{d_{z,\lambda}+1}),
\end{equation}
where $d_{z,\lambda}\in\{0,1,2,\ldots\}$ and
$\eta_{z,\lambda}>0$. If $\theta_{z,\lambda}>0$, then
$d_{z,\lambda}=0$ and $\eta_{z,\lambda}=\theta_{z,\lambda}$.
If $\theta_{z,\lambda}=0$ and the nonlinear germ is nonzero, then
$d_{z,\lambda}\ge1$. Every such nonlinear germ is retained.
\end{theorem}
\begin{proof}
Insert the fixed angular cut as two constant endpoints in the ordered
list of Lemma~\ref{lem:v68-root-order}. The angular fraction of each
label is the sum, divided by $2\pi$, of the lengths of those open
arcs carrying that label. Each length is a difference of analytic
root lifts, or of a root lift and a constant endpoint. Hence this
finite sum is analytic at zero. Every summand is nonnegative for
positive $r$, including arcs whose endpoints coincide at zero. Its
constant term is the tangent angular fraction: away from the finite
set of tangent zero directions, all signs converge to their tangent
signs, with the inactive signs fixed. This also follows directly by
dominated convergence in \eqref{eq:v68-unweighted-angular-germ}.

An analytic scalar germ is either identically zero or has a first
nonzero Taylor coefficient. Since the angular fraction is nonnegative
for positive $r$, that coefficient is positive. These observations
prove all assertions. In the zero-germ case the corresponding physical
angular source is zero on this smaller disk; no assertion is made
about the rest of the original chart.
\end{proof}

\subsection{A recorded radius and a relative Taylor remainder}
For a nonzero germ fix $0<r_{z,\lambda}^*<r_\chi$ within all the
preceding analytic, ordering and inactive-sign radii. Put $d=d_{z,\lambda}$,
$\eta=\eta_{z,\lambda}$ and
\begin{equation}\label{eq:v68-angular-relative-budget}
 L_{z,\lambda}=\frac{1}{(d+1)!}
       \sup_{0\le r\le r_{z,\lambda}^*}
                       |\ell_{z,\lambda}^{(d+1)}(r)|,
 \qquad k_{z,\lambda}^{\rm a}=L_{z,\lambda}/\eta.
\end{equation}
Taylor's theorem, with the lower coefficients zero, gives
\begin{equation}\label{eq:v68-relative-angular-type}
 |\ell_{z,\lambda}(r)-\eta r^d|
          \le k_{z,\lambda}^{\rm a}\,r\,\eta r^d
                  \quad(0<r<r_{z,\lambda}^*).
\end{equation}
This is relative to the first nonzero coefficient, not to a vanished
tangent fraction.

The data have a finite certificate at each nonzero germ: the primitive
analytic formulas, the chosen radius, the cyclic ordering of the root
germs, their first separating jets, the positive coefficient $\eta$,
and the derivative remainder in \eqref{eq:v68-angular-relative-budget}.
An asserted equality of two root germs requires an analytic identity;
it is not certified by an arbitrary finite list of vanishing derivatives.
For a nonidentical pair, its first separating coefficient and a remainder
bound certify the order on a sufficiently small radius. Root derivatives
are obtained by differentiating $F_f(r,\phi_f(r))=0$; the divisor at
zero is the displayed nonzero angular derivative. No lower bound on
that divisor, no uniform maximum separating order, and no uniform
radius is hidden in the certificate. The radius and remainder are
recorded separately and will enter the physical cutoff.

\begin{remark}[Arbitrarily high birth order is locally possible]
\label{rem:v68-high-order-birth}
The analytic sign conditions $x>0$, $y>0$, and $y<x^{d+1}$ have
nonzero active gradients for every integer $d\ge1$. On the circle
of radius $r$, their angular width is the small positive solution
$\alpha(r)$ of
$\sin\alpha=r^d\cos^{d+1}\alpha$. Thus
$\alpha(r)/(2\pi)=r^d/(2\pi)+O(r^{d+1})$ and its tangent fraction
is zero. This local example illustrates the classification; it is not
claimed to construct a Lorentz itinerary of every possible order.
The following height theorem does not require such a realization or
a bound on the order.
\end{remark}
